% appendix_schema_enum_mapping.tex
% Appendix: JSON Schema Enum Mapping Table
% Addresses Locked Decision #7: "Do NOT rename JSON schema enums. Add explicit
% mapping table in paper Appendix."
%
% Include with: % appendix_schema_enum_mapping.tex
% Appendix: JSON Schema Enum Mapping Table
% Addresses Locked Decision #7: "Do NOT rename JSON schema enums. Add explicit
% mapping table in paper Appendix."
%
% Include with: % appendix_schema_enum_mapping.tex
% Appendix: JSON Schema Enum Mapping Table
% Addresses Locked Decision #7: "Do NOT rename JSON schema enums. Add explicit
% mapping table in paper Appendix."
%
% Include with: % appendix_schema_enum_mapping.tex
% Appendix: JSON Schema Enum Mapping Table
% Addresses Locked Decision #7: "Do NOT rename JSON schema enums. Add explicit
% mapping table in paper Appendix."
%
% Include with: \input{appendix_schema_enum_mapping} inside \appendix block
% Cross-reference: \ref{sec:appendix-schema-enums}

\section{JSON Schema Enum Mapping Reference}
\label{sec:appendix-schema-enums}

\noindent
The IPIBench dataset (\texttt{benchmark.json}) uses compact machine-readable enum
values across all ten taxonomy dimensions.
To support reproducibility and prevent ambiguity between the paper's prose
descriptions and the programmatic values, this appendix provides the authoritative
mapping from every schema enum to its human-readable paper label.
\textbf{These enum strings are frozen as of Benchmark v1.0 and must not be renamed
or aliased without incrementing the dataset version.}

% ------------------------------------------------------------------
\subsection*{A.1 \texttt{category} -- Content Delivery Vector}
% ------------------------------------------------------------------

\begin{table}[h]
\centering
\caption{Enum mapping for \texttt{category} (content delivery vector).}
\label{tab:enum-category}
\begin{tabular}{@{}ll@{}}
\toprule
\textbf{Schema Enum Value} & \textbf{Paper Label} \\
\midrule
\texttt{webpage}      & Web Document Injection \\
\texttt{file}         & File Content Injection \\
\texttt{tool\_output} & Tool / API Output Injection \\
\bottomrule
\end{tabular}
\end{table}

% ------------------------------------------------------------------
\subsection*{A.2 \texttt{attack\_goal} -- Primary Attacker Objective}
% ------------------------------------------------------------------

\begin{table}[h]
\centering
\caption{Enum mapping for \texttt{attack\_goal} (primary attacker objective).}
\label{tab:enum-attack-goal}
\begin{tabular}{@{}lp{8cm}@{}}
\toprule
\textbf{Schema Enum Value} & \textbf{Paper Label \& Description} \\
\midrule
\texttt{task\_hijack}          & Task Hijacking -- redirect the model's primary task to an attacker-specified action (64 scenarios) \\
\texttt{data\_exfiltration}    & Information Disclosure Compliance -- elicit disclosure of system prompt, session context, or user data (18 scenarios) \\
\texttt{unauthorized\_action}  & Unauthorized Action Induction -- trigger a simulated tool call or side-effect not authorized by the user (9 scenarios) \\
\texttt{persistence}           & Stateful Persistence Injection -- embed instructions intended to corrupt future conversational turns (4 scenarios) \\
\texttt{identity\_corruption}  & Identity Corruption -- overwrite the model's stated identity, persona, or capability claims (5 scenarios; 2 held out) \\
\bottomrule
\end{tabular}
\end{table}

% ------------------------------------------------------------------
\subsection*{A.3 \texttt{evasion\_style} -- Injection Evasion Technique}
% ------------------------------------------------------------------

\begin{table}[h]
\centering
\caption{Enum mapping for \texttt{evasion\_style} (injection evasion technique).}
\label{tab:enum-evasion-style}
\begin{tabular}{@{}lp{8cm}@{}}
\toprule
\textbf{Schema Enum Value} & \textbf{Paper Label \& Description} \\
\midrule
\texttt{direct}                   & Direct Instruction -- plaintext instruction with no disguise \\
\texttt{authority\_impersonation} & Authority Impersonation -- injection claims to originate from a trusted entity (vendor, safety team) \\
\texttt{role\_mimicry}            & Role Mimicry -- injection instructs the model to adopt an alternate persona or role \\
\texttt{hypothetical\_framing}    & Hypothetical Framing -- embeds the malicious instruction inside a fictional or counterfactual premise \\
\texttt{obfuscated}               & Obfuscated Encoding -- uses Base64, leetspeak, Unicode homoglyphs, or other encoding to bypass lexical filters \\
\texttt{multilingual}             & Cross-Lingual Evasion -- injection delivered in a non-English language to evade English-dominant safety filters \\
\bottomrule
\end{tabular}
\end{table}

% ------------------------------------------------------------------
\subsection*{A.4 \texttt{defense\_mode} -- Evaluation Defense Configuration}
% ------------------------------------------------------------------

\begin{table}[h]
\centering
\caption{Enum mapping for \texttt{defense\_mode} (evaluation defense configuration).}
\label{tab:enum-defense-mode}
\begin{tabular}{@{}lp{8.5cm}@{}}
\toprule
\textbf{Schema Enum Value} & \textbf{Paper Label \& Description} \\
\midrule
\texttt{none}            & Baseline (No Defense) -- empty system prompt, unmodified context \\
\texttt{prompt\_warning} & Prompt Warning -- system prompt includes an explicit security notice; no structural modification \\
\texttt{spotlighting}    & Spotlighting Delimiter -- external content wrapped in \texttt{<EXTERNAL\_DATA>} XML delimiters; system prompt instructs the model to ignore commands inside delimiters \\
\texttt{input\_filter}   & Input Filter -- known injection trigger phrases replaced with \texttt{[FILTERED]} before delivery to the model \\
\bottomrule
\end{tabular}
\end{table}

% ------------------------------------------------------------------
\subsection*{A.5 \texttt{injection\_position} -- Position of Injected Payload}
% ------------------------------------------------------------------

\begin{table}[h]
\centering
\caption{Enum mapping for \texttt{injection\_position}.}
\label{tab:enum-injection-position}
\begin{tabular}{@{}ll@{}}
\toprule
\textbf{Schema Enum Value} & \textbf{Paper Label} \\
\midrule
\texttt{prefix} & Prefix -- injection appears before legitimate content \\
\texttt{suffix} & Suffix -- injection appended after legitimate content \\
\texttt{inline} & Inline -- injection embedded mid-document, mixed with legitimate content \\
\bottomrule
\end{tabular}
\end{table}

% ------------------------------------------------------------------
\subsection*{A.6 Secondary Metadata Dimensions}
% ------------------------------------------------------------------

\noindent
The following dimensions exhibit lower empirical variance ($>65\%$ modal class
concentration) and are retained for schema filtering only, not primary hypothesis
testing (see Limitations, \S\ref{sec:lim-taxonomy}).

\begin{table}[h]
\centering
\caption{Secondary enum dimensions: authority\_claimed, target\_action\_type, linguistic\_register, harm\_severity, persistence.}
\label{tab:enum-secondary}
\begin{tabular}{@{}lll@{}}
\toprule
\textbf{Dimension} & \textbf{Schema Enum Values} & \textbf{Notes} \\
\midrule
\texttt{authority\_claimed}   & \texttt{none}, \texttt{vendor}, \texttt{system}, \texttt{user}                                                        & Whether injection claims institutional authority \\
\texttt{target\_action\_type} & \texttt{output\_modification}, \texttt{tool\_invocation}, \texttt{data\_disclosure}, \texttt{session\_corruption} & Mechanical action type targeted by injection \\
\texttt{linguistic\_register} & \texttt{formal}, \texttt{casual}, \texttt{technical}, \texttt{urgent}                                                & Stylistic register of the injection text \\
\texttt{harm\_severity}       & \texttt{low}, \texttt{medium}, \texttt{high}                                                                          & Potential downstream harm if injection succeeds \\
\texttt{persistence}          & \texttt{true}, \texttt{false}                                                                                         & Whether the injection persists across turns \\
\bottomrule
\end{tabular}
\end{table}
 inside \appendix block
% Cross-reference: \ref{sec:appendix-schema-enums}

\section{JSON Schema Enum Mapping Reference}
\label{sec:appendix-schema-enums}

\noindent
The IPIBench dataset (\texttt{benchmark.json}) uses compact machine-readable enum
values across all ten taxonomy dimensions.
To support reproducibility and prevent ambiguity between the paper's prose
descriptions and the programmatic values, this appendix provides the authoritative
mapping from every schema enum to its human-readable paper label.
\textbf{These enum strings are frozen as of Benchmark v1.0 and must not be renamed
or aliased without incrementing the dataset version.}

% ------------------------------------------------------------------
\subsection*{A.1 \texttt{category} -- Content Delivery Vector}
% ------------------------------------------------------------------

\begin{table}[h]
\centering
\caption{Enum mapping for \texttt{category} (content delivery vector).}
\label{tab:enum-category}
\begin{tabular}{@{}ll@{}}
\toprule
\textbf{Schema Enum Value} & \textbf{Paper Label} \\
\midrule
\texttt{webpage}      & Web Document Injection \\
\texttt{file}         & File Content Injection \\
\texttt{tool\_output} & Tool / API Output Injection \\
\bottomrule
\end{tabular}
\end{table}

% ------------------------------------------------------------------
\subsection*{A.2 \texttt{attack\_goal} -- Primary Attacker Objective}
% ------------------------------------------------------------------

\begin{table}[h]
\centering
\caption{Enum mapping for \texttt{attack\_goal} (primary attacker objective).}
\label{tab:enum-attack-goal}
\begin{tabular}{@{}lp{8cm}@{}}
\toprule
\textbf{Schema Enum Value} & \textbf{Paper Label \& Description} \\
\midrule
\texttt{task\_hijack}          & Task Hijacking -- redirect the model's primary task to an attacker-specified action (64 scenarios) \\
\texttt{data\_exfiltration}    & Information Disclosure Compliance -- elicit disclosure of system prompt, session context, or user data (18 scenarios) \\
\texttt{unauthorized\_action}  & Unauthorized Action Induction -- trigger a simulated tool call or side-effect not authorized by the user (9 scenarios) \\
\texttt{persistence}           & Stateful Persistence Injection -- embed instructions intended to corrupt future conversational turns (4 scenarios) \\
\texttt{identity\_corruption}  & Identity Corruption -- overwrite the model's stated identity, persona, or capability claims (5 scenarios; 2 held out) \\
\bottomrule
\end{tabular}
\end{table}

% ------------------------------------------------------------------
\subsection*{A.3 \texttt{evasion\_style} -- Injection Evasion Technique}
% ------------------------------------------------------------------

\begin{table}[h]
\centering
\caption{Enum mapping for \texttt{evasion\_style} (injection evasion technique).}
\label{tab:enum-evasion-style}
\begin{tabular}{@{}lp{8cm}@{}}
\toprule
\textbf{Schema Enum Value} & \textbf{Paper Label \& Description} \\
\midrule
\texttt{direct}                   & Direct Instruction -- plaintext instruction with no disguise \\
\texttt{authority\_impersonation} & Authority Impersonation -- injection claims to originate from a trusted entity (vendor, safety team) \\
\texttt{role\_mimicry}            & Role Mimicry -- injection instructs the model to adopt an alternate persona or role \\
\texttt{hypothetical\_framing}    & Hypothetical Framing -- embeds the malicious instruction inside a fictional or counterfactual premise \\
\texttt{obfuscated}               & Obfuscated Encoding -- uses Base64, leetspeak, Unicode homoglyphs, or other encoding to bypass lexical filters \\
\texttt{multilingual}             & Cross-Lingual Evasion -- injection delivered in a non-English language to evade English-dominant safety filters \\
\bottomrule
\end{tabular}
\end{table}

% ------------------------------------------------------------------
\subsection*{A.4 \texttt{defense\_mode} -- Evaluation Defense Configuration}
% ------------------------------------------------------------------

\begin{table}[h]
\centering
\caption{Enum mapping for \texttt{defense\_mode} (evaluation defense configuration).}
\label{tab:enum-defense-mode}
\begin{tabular}{@{}lp{8.5cm}@{}}
\toprule
\textbf{Schema Enum Value} & \textbf{Paper Label \& Description} \\
\midrule
\texttt{none}            & Baseline (No Defense) -- empty system prompt, unmodified context \\
\texttt{prompt\_warning} & Prompt Warning -- system prompt includes an explicit security notice; no structural modification \\
\texttt{spotlighting}    & Spotlighting Delimiter -- external content wrapped in \texttt{<EXTERNAL\_DATA>} XML delimiters; system prompt instructs the model to ignore commands inside delimiters \\
\texttt{input\_filter}   & Input Filter -- known injection trigger phrases replaced with \texttt{[FILTERED]} before delivery to the model \\
\bottomrule
\end{tabular}
\end{table}

% ------------------------------------------------------------------
\subsection*{A.5 \texttt{injection\_position} -- Position of Injected Payload}
% ------------------------------------------------------------------

\begin{table}[h]
\centering
\caption{Enum mapping for \texttt{injection\_position}.}
\label{tab:enum-injection-position}
\begin{tabular}{@{}ll@{}}
\toprule
\textbf{Schema Enum Value} & \textbf{Paper Label} \\
\midrule
\texttt{prefix} & Prefix -- injection appears before legitimate content \\
\texttt{suffix} & Suffix -- injection appended after legitimate content \\
\texttt{inline} & Inline -- injection embedded mid-document, mixed with legitimate content \\
\bottomrule
\end{tabular}
\end{table}

% ------------------------------------------------------------------
\subsection*{A.6 Secondary Metadata Dimensions}
% ------------------------------------------------------------------

\noindent
The following dimensions exhibit lower empirical variance ($>65\%$ modal class
concentration) and are retained for schema filtering only, not primary hypothesis
testing (see Limitations, \S\ref{sec:lim-taxonomy}).

\begin{table}[h]
\centering
\caption{Secondary enum dimensions: authority\_claimed, target\_action\_type, linguistic\_register, harm\_severity, persistence.}
\label{tab:enum-secondary}
\begin{tabular}{@{}lll@{}}
\toprule
\textbf{Dimension} & \textbf{Schema Enum Values} & \textbf{Notes} \\
\midrule
\texttt{authority\_claimed}   & \texttt{none}, \texttt{vendor}, \texttt{system}, \texttt{user}                                                        & Whether injection claims institutional authority \\
\texttt{target\_action\_type} & \texttt{output\_modification}, \texttt{tool\_invocation}, \texttt{data\_disclosure}, \texttt{session\_corruption} & Mechanical action type targeted by injection \\
\texttt{linguistic\_register} & \texttt{formal}, \texttt{casual}, \texttt{technical}, \texttt{urgent}                                                & Stylistic register of the injection text \\
\texttt{harm\_severity}       & \texttt{low}, \texttt{medium}, \texttt{high}                                                                          & Potential downstream harm if injection succeeds \\
\texttt{persistence}          & \texttt{true}, \texttt{false}                                                                                         & Whether the injection persists across turns \\
\bottomrule
\end{tabular}
\end{table}
 inside \appendix block
% Cross-reference: \ref{sec:appendix-schema-enums}

\section{JSON Schema Enum Mapping Reference}
\label{sec:appendix-schema-enums}

\noindent
The IPIBench dataset (\texttt{benchmark.json}) uses compact machine-readable enum
values across all ten taxonomy dimensions.
To support reproducibility and prevent ambiguity between the paper's prose
descriptions and the programmatic values, this appendix provides the authoritative
mapping from every schema enum to its human-readable paper label.
\textbf{These enum strings are frozen as of Benchmark v1.0 and must not be renamed
or aliased without incrementing the dataset version.}

% ------------------------------------------------------------------
\subsection*{A.1 \texttt{category} -- Content Delivery Vector}
% ------------------------------------------------------------------

\begin{table}[h]
\centering
\caption{Enum mapping for \texttt{category} (content delivery vector).}
\label{tab:enum-category}
\begin{tabular}{@{}ll@{}}
\toprule
\textbf{Schema Enum Value} & \textbf{Paper Label} \\
\midrule
\texttt{webpage}      & Web Document Injection \\
\texttt{file}         & File Content Injection \\
\texttt{tool\_output} & Tool / API Output Injection \\
\bottomrule
\end{tabular}
\end{table}

% ------------------------------------------------------------------
\subsection*{A.2 \texttt{attack\_goal} -- Primary Attacker Objective}
% ------------------------------------------------------------------

\begin{table}[h]
\centering
\caption{Enum mapping for \texttt{attack\_goal} (primary attacker objective).}
\label{tab:enum-attack-goal}
\begin{tabular}{@{}lp{8cm}@{}}
\toprule
\textbf{Schema Enum Value} & \textbf{Paper Label \& Description} \\
\midrule
\texttt{task\_hijack}          & Task Hijacking -- redirect the model's primary task to an attacker-specified action (64 scenarios) \\
\texttt{data\_exfiltration}    & Information Disclosure Compliance -- elicit disclosure of system prompt, session context, or user data (18 scenarios) \\
\texttt{unauthorized\_action}  & Unauthorized Action Induction -- trigger a simulated tool call or side-effect not authorized by the user (9 scenarios) \\
\texttt{persistence}           & Stateful Persistence Injection -- embed instructions intended to corrupt future conversational turns (4 scenarios) \\
\texttt{identity\_corruption}  & Identity Corruption -- overwrite the model's stated identity, persona, or capability claims (5 scenarios; 2 held out) \\
\bottomrule
\end{tabular}
\end{table}

% ------------------------------------------------------------------
\subsection*{A.3 \texttt{evasion\_style} -- Injection Evasion Technique}
% ------------------------------------------------------------------

\begin{table}[h]
\centering
\caption{Enum mapping for \texttt{evasion\_style} (injection evasion technique).}
\label{tab:enum-evasion-style}
\begin{tabular}{@{}lp{8cm}@{}}
\toprule
\textbf{Schema Enum Value} & \textbf{Paper Label \& Description} \\
\midrule
\texttt{direct}                   & Direct Instruction -- plaintext instruction with no disguise \\
\texttt{authority\_impersonation} & Authority Impersonation -- injection claims to originate from a trusted entity (vendor, safety team) \\
\texttt{role\_mimicry}            & Role Mimicry -- injection instructs the model to adopt an alternate persona or role \\
\texttt{hypothetical\_framing}    & Hypothetical Framing -- embeds the malicious instruction inside a fictional or counterfactual premise \\
\texttt{obfuscated}               & Obfuscated Encoding -- uses Base64, leetspeak, Unicode homoglyphs, or other encoding to bypass lexical filters \\
\texttt{multilingual}             & Cross-Lingual Evasion -- injection delivered in a non-English language to evade English-dominant safety filters \\
\bottomrule
\end{tabular}
\end{table}

% ------------------------------------------------------------------
\subsection*{A.4 \texttt{defense\_mode} -- Evaluation Defense Configuration}
% ------------------------------------------------------------------

\begin{table}[h]
\centering
\caption{Enum mapping for \texttt{defense\_mode} (evaluation defense configuration).}
\label{tab:enum-defense-mode}
\begin{tabular}{@{}lp{8.5cm}@{}}
\toprule
\textbf{Schema Enum Value} & \textbf{Paper Label \& Description} \\
\midrule
\texttt{none}            & Baseline (No Defense) -- empty system prompt, unmodified context \\
\texttt{prompt\_warning} & Prompt Warning -- system prompt includes an explicit security notice; no structural modification \\
\texttt{spotlighting}    & Spotlighting Delimiter -- external content wrapped in \texttt{<EXTERNAL\_DATA>} XML delimiters; system prompt instructs the model to ignore commands inside delimiters \\
\texttt{input\_filter}   & Input Filter -- known injection trigger phrases replaced with \texttt{[FILTERED]} before delivery to the model \\
\bottomrule
\end{tabular}
\end{table}

% ------------------------------------------------------------------
\subsection*{A.5 \texttt{injection\_position} -- Position of Injected Payload}
% ------------------------------------------------------------------

\begin{table}[h]
\centering
\caption{Enum mapping for \texttt{injection\_position}.}
\label{tab:enum-injection-position}
\begin{tabular}{@{}ll@{}}
\toprule
\textbf{Schema Enum Value} & \textbf{Paper Label} \\
\midrule
\texttt{prefix} & Prefix -- injection appears before legitimate content \\
\texttt{suffix} & Suffix -- injection appended after legitimate content \\
\texttt{inline} & Inline -- injection embedded mid-document, mixed with legitimate content \\
\bottomrule
\end{tabular}
\end{table}

% ------------------------------------------------------------------
\subsection*{A.6 Secondary Metadata Dimensions}
% ------------------------------------------------------------------

\noindent
The following dimensions exhibit lower empirical variance ($>65\%$ modal class
concentration) and are retained for schema filtering only, not primary hypothesis
testing (see Limitations, \S\ref{sec:lim-taxonomy}).

\begin{table}[h]
\centering
\caption{Secondary enum dimensions: authority\_claimed, target\_action\_type, linguistic\_register, harm\_severity, persistence.}
\label{tab:enum-secondary}
\begin{tabular}{@{}lll@{}}
\toprule
\textbf{Dimension} & \textbf{Schema Enum Values} & \textbf{Notes} \\
\midrule
\texttt{authority\_claimed}   & \texttt{none}, \texttt{vendor}, \texttt{system}, \texttt{user}                                                        & Whether injection claims institutional authority \\
\texttt{target\_action\_type} & \texttt{output\_modification}, \texttt{tool\_invocation}, \texttt{data\_disclosure}, \texttt{session\_corruption} & Mechanical action type targeted by injection \\
\texttt{linguistic\_register} & \texttt{formal}, \texttt{casual}, \texttt{technical}, \texttt{urgent}                                                & Stylistic register of the injection text \\
\texttt{harm\_severity}       & \texttt{low}, \texttt{medium}, \texttt{high}                                                                          & Potential downstream harm if injection succeeds \\
\texttt{persistence}          & \texttt{true}, \texttt{false}                                                                                         & Whether the injection persists across turns \\
\bottomrule
\end{tabular}
\end{table}
 inside \appendix block
% Cross-reference: \ref{sec:appendix-schema-enums}

\section{JSON Schema Enum Mapping Reference}
\label{sec:appendix-schema-enums}

\noindent
The IPIBench dataset (\texttt{benchmark.json}) uses compact machine-readable enum
values across all ten taxonomy dimensions.
To support reproducibility and prevent ambiguity between the paper's prose
descriptions and the programmatic values, this appendix provides the authoritative
mapping from every schema enum to its human-readable paper label.
\textbf{These enum strings are frozen as of Benchmark v1.0 and must not be renamed
or aliased without incrementing the dataset version.}

% ------------------------------------------------------------------
\subsection*{A.1 \texttt{category} -- Content Delivery Vector}
% ------------------------------------------------------------------

\begin{table}[h]
\centering
\caption{Enum mapping for \texttt{category} (content delivery vector).}
\label{tab:enum-category}
\begin{tabular}{@{}ll@{}}
\toprule
\textbf{Schema Enum Value} & \textbf{Paper Label} \\
\midrule
\texttt{webpage}      & Web Document Injection \\
\texttt{file}         & File Content Injection \\
\texttt{tool\_output} & Tool / API Output Injection \\
\bottomrule
\end{tabular}
\end{table}

% ------------------------------------------------------------------
\subsection*{A.2 \texttt{attack\_goal} -- Primary Attacker Objective}
% ------------------------------------------------------------------

\begin{table}[h]
\centering
\caption{Enum mapping for \texttt{attack\_goal} (primary attacker objective).}
\label{tab:enum-attack-goal}
\begin{tabular}{@{}lp{8cm}@{}}
\toprule
\textbf{Schema Enum Value} & \textbf{Paper Label \& Description} \\
\midrule
\texttt{task\_hijack}          & Task Hijacking -- redirect the model's primary task to an attacker-specified action (64 scenarios) \\
\texttt{data\_exfiltration}    & Information Disclosure Compliance -- elicit disclosure of system prompt, session context, or user data (18 scenarios) \\
\texttt{unauthorized\_action}  & Unauthorized Action Induction -- trigger a simulated tool call or side-effect not authorized by the user (9 scenarios) \\
\texttt{persistence}           & Stateful Persistence Injection -- embed instructions intended to corrupt future conversational turns (4 scenarios) \\
\texttt{identity\_corruption}  & Identity Corruption -- overwrite the model's stated identity, persona, or capability claims (5 scenarios; 2 held out) \\
\bottomrule
\end{tabular}
\end{table}

% ------------------------------------------------------------------
\subsection*{A.3 \texttt{evasion\_style} -- Injection Evasion Technique}
% ------------------------------------------------------------------

\begin{table}[h]
\centering
\caption{Enum mapping for \texttt{evasion\_style} (injection evasion technique).}
\label{tab:enum-evasion-style}
\begin{tabular}{@{}lp{8cm}@{}}
\toprule
\textbf{Schema Enum Value} & \textbf{Paper Label \& Description} \\
\midrule
\texttt{direct}                   & Direct Instruction -- plaintext instruction with no disguise \\
\texttt{authority\_impersonation} & Authority Impersonation -- injection claims to originate from a trusted entity (vendor, safety team) \\
\texttt{role\_mimicry}            & Role Mimicry -- injection instructs the model to adopt an alternate persona or role \\
\texttt{hypothetical\_framing}    & Hypothetical Framing -- embeds the malicious instruction inside a fictional or counterfactual premise \\
\texttt{obfuscated}               & Obfuscated Encoding -- uses Base64, leetspeak, Unicode homoglyphs, or other encoding to bypass lexical filters \\
\texttt{multilingual}             & Cross-Lingual Evasion -- injection delivered in a non-English language to evade English-dominant safety filters \\
\bottomrule
\end{tabular}
\end{table}

% ------------------------------------------------------------------
\subsection*{A.4 \texttt{defense\_mode} -- Evaluation Defense Configuration}
% ------------------------------------------------------------------

\begin{table}[h]
\centering
\caption{Enum mapping for \texttt{defense\_mode} (evaluation defense configuration).}
\label{tab:enum-defense-mode}
\begin{tabular}{@{}lp{8.5cm}@{}}
\toprule
\textbf{Schema Enum Value} & \textbf{Paper Label \& Description} \\
\midrule
\texttt{none}            & Baseline (No Defense) -- empty system prompt, unmodified context \\
\texttt{prompt\_warning} & Prompt Warning -- system prompt includes an explicit security notice; no structural modification \\
\texttt{spotlighting}    & Spotlighting Delimiter -- external content wrapped in \texttt{<EXTERNAL\_DATA>} XML delimiters; system prompt instructs the model to ignore commands inside delimiters \\
\texttt{input\_filter}   & Input Filter -- known injection trigger phrases replaced with \texttt{[FILTERED]} before delivery to the model \\
\bottomrule
\end{tabular}
\end{table}

% ------------------------------------------------------------------
\subsection*{A.5 \texttt{injection\_position} -- Position of Injected Payload}
% ------------------------------------------------------------------

\begin{table}[h]
\centering
\caption{Enum mapping for \texttt{injection\_position}.}
\label{tab:enum-injection-position}
\begin{tabular}{@{}ll@{}}
\toprule
\textbf{Schema Enum Value} & \textbf{Paper Label} \\
\midrule
\texttt{prefix} & Prefix -- injection appears before legitimate content \\
\texttt{suffix} & Suffix -- injection appended after legitimate content \\
\texttt{inline} & Inline -- injection embedded mid-document, mixed with legitimate content \\
\bottomrule
\end{tabular}
\end{table}

% ------------------------------------------------------------------
\subsection*{A.6 Secondary Metadata Dimensions}
% ------------------------------------------------------------------

\noindent
The following dimensions exhibit lower empirical variance ($>65\%$ modal class
concentration) and are retained for schema filtering only, not primary hypothesis
testing (see Limitations, \S\ref{sec:lim-taxonomy}).

\begin{table}[h]
\centering
\caption{Secondary enum dimensions: authority\_claimed, target\_action\_type, linguistic\_register, harm\_severity, persistence.}
\label{tab:enum-secondary}
\begin{tabular}{@{}lll@{}}
\toprule
\textbf{Dimension} & \textbf{Schema Enum Values} & \textbf{Notes} \\
\midrule
\texttt{authority\_claimed}   & \texttt{none}, \texttt{vendor}, \texttt{system}, \texttt{user}                                                        & Whether injection claims institutional authority \\
\texttt{target\_action\_type} & \texttt{output\_modification}, \texttt{tool\_invocation}, \texttt{data\_disclosure}, \texttt{session\_corruption} & Mechanical action type targeted by injection \\
\texttt{linguistic\_register} & \texttt{formal}, \texttt{casual}, \texttt{technical}, \texttt{urgent}                                                & Stylistic register of the injection text \\
\texttt{harm\_severity}       & \texttt{low}, \texttt{medium}, \texttt{high}                                                                          & Potential downstream harm if injection succeeds \\
\texttt{persistence}          & \texttt{true}, \texttt{false}                                                                                         & Whether the injection persists across turns \\
\bottomrule
\end{tabular}
\end{table}
